\section{Uniform convergence}
\label{sec:uniform-convergence-new}
\textbf{Before starting.} Pointwise convergence means fixing one input $x$
first and then letting $n$ grow. Uniform convergence means finding one index
that works for \emph{all} inputs at once. We shall repeatedly estimate the
largest error
\[
 E_n=\sup_{x\in S}|f_n(x)-f(x)|.
\]
Uniform convergence is equivalent to $E_n\to0$; Exercise~24.12 proves this.
To disprove uniform convergence, it is enough to find a positive number
$\varepsilon_0$ and inputs $x_n\in S$, for arbitrarily large $n$, with
$|f_n(x_n)-f(x_n)|\ge\varepsilon_0$. The input is allowed to change with $n$.
That is precisely the obstacle uniform convergence must overcome.

\exercise{24.1}{One bound for every oscillation}
\goal Prove uniform convergence to zero for
$f_n(x)=(1+2\cos^2(nx))/\sqrt n$ on $\R$.
\idea We need not follow the oscillations: the cosine is always between $-1$ and $1$.
\step{Bound the numerator independently of the input.}
Since $0\le\cos^2(nx)\le1$, we have
\[
 |f_n(x)-0|\le\frac3{\sqrt n}\qquad(x\in\R).
\]
Given $\varepsilon>0$, choose an integer $N>(3/\varepsilon)^2$.
For every $n>N$ and every real $x$, the last bound is smaller than
$\varepsilon$. The choice of $N$ depends only on $\varepsilon$.
\answer{$f_n\to0$ uniformly on $\R$.}

\exercise{24.2}{A small slope on bounded and unbounded domains}
\goal Analyze $f_n(x)=x/n$ on $[0,\infty)$.
\step{(a) Fix the input before passing to the limit.}
For fixed $x\ge0$, the numerator is constant and $n\to\infty$, so
$f_n(x)\to0$. Thus the pointwise limit is $f(x)=0$.
\step{(b) Control all inputs in $[0,1]$.}
There $|f_n(x)|\le1/n$. Given $\varepsilon>0$, any integer
$N>1/\varepsilon$ works for all $x\in[0,1]$ and $n>N$.
\step{(c) Let the input escape along with the index.}
On $[0,\infty)$ take $x_n=n$. Then $f_n(x_n)=1$, so the error never
falls below $1/2$. This contradicts the requirement of uniform convergence.
\answer{The limit is zero everywhere; convergence is uniform on $[0,1]$
but not on $[0,\infty)$.}

\exercise{24.3}{A transition concentrated near one}
\goal Analyze $f_n(x)=1/(1+x^n)$ for $x\ge0$.
\step{(a) Separate the three possible behaviors of a power.}
If $0\le x<1$, then $x^n\to0$, giving limit $1$.
At $x=1$ every term is $1/2$.
If $x>1$, then $x^n\to\infty$, giving limit $0$. Thus
\[
 f(x)=\begin{cases}1&0\le x<1,\\1/2&x=1,\\0&x>1.\end{cases}
\]
\step{(b) Test uniform convergence on $[0,1]$.}
Each $f_n$ is continuous there because its denominator is positive.
The limit is not continuous at $1$: its values from the left are $1$,
whereas $f(1)=1/2$. The uniform-limit theorem for continuous functions
therefore excludes uniform convergence on $[0,1]$.
\step{(c) Use restriction to a smaller domain.}
Uniform convergence on $[0,\infty)$ would imply uniform convergence on
its subset $[0,1]$, with exactly the same error bound. Part (b) rules that out.
\answer{The piecewise function above is the pointwise limit, and neither
of the two requested uniform-convergence assertions holds.}

\exercise{24.4}{The complementary transition}
\goal Analyze $f_n(x)=x^n/(1+x^n)$ on $[0,\infty)$.
\step{(a) Use an exact identity.}
We have $f_n(x)=1-1/(1+x^n)$. The three cases of the preceding
exercise therefore give
\[
 f(x)=\begin{cases}0&0\le x<1,\\1/2&x=1,\\1&x>1.\end{cases}
\]
\step{(b) Notice the endpoint jump.}
The $f_n$ are continuous on $[0,1]$, but $f$ is not continuous at $1$.
Consequently convergence cannot be uniform there, by Theorem~24.3.
\step{(c) Restrict any proposed larger-domain convergence.}
Uniform convergence on $[0,\infty)$ would restrict to uniform convergence
on $[0,1]$, contradicting part (b).
\answer{The displayed $f$ is the limit; convergence is uniform on neither
$[0,1]$ nor $[0,\infty)$.}

\exercise{24.5}{A denominator containing both the index and a power}
\goal Analyze $f_n(x)=x^n/(n+x^n)$ for $x\ge0$.
\step{(a) Find the pointwise limit.}
For $0\le x\le1$, we have $0\le f_n(x)\le1/n$, so the limit is zero.
For fixed $x>1$, divide numerator and denominator by $x^n$:
\[
 f_n(x)=\frac1{1+n/x^n}\longrightarrow1.
\]
Here $n/x^n\to0$ because exponential growth dominates the linear factor;
for example, the ratio of successive terms tends to $1/x<1$.
Thus $f(x)=0$ on $[0,1]$ and $f(x)=1$ on $(1,\infty)$.
\step{(b) Obtain a common bound on $[0,1]$.}
The estimate $|f_n(x)-f(x)|\le1/n$ works simultaneously for every input
in that interval. Hence convergence is uniform there.
\step{(c) Use the jump at an interior point of the larger domain.}
Every $f_n$ is continuous on $[0,\infty)$, but the pointwise limit jumps
at $1$. The uniform-limit theorem excludes uniform convergence on the larger domain.
\answer{Uniform convergence holds on $[0,1]$, and fails on $[0,\infty)$.}

\exercise{24.6}{Shifted squares}
\goal Find the limit of $f_n(x)=(x-1/n)^2$ on $[0,1]$ and decide uniform convergence.
\step{(a) Expand to identify the limit.}
For fixed $x$, $f_n(x)=x^2-2x/n+1/n^2\to x^2$. Thus $f(x)=x^2$.
\step{(b) Bound the entire error, rather than just each separate input.}
For $0\le x\le1$,
\[
 |f_n(x)-x^2|=\left|-\frac{2x}{n}+\frac1{n^2}\right|
 \le\frac2n+\frac1{n^2}\le\frac3n.
\]
Given $\varepsilon>0$, choose $N>3/\varepsilon$. This bound is then
less than $\varepsilon$ for every $x\in[0,1]$ and $n>N$.
\answer{$f_n\to(x\mapsto x^2)$ uniformly on $[0,1]$.}

\exercise{24.7}{Subtracting a power}
\goal Analyze $f_n(x)=x-x^n$ on $[0,1]$.
\step{(a) Check the endpoint separately.}
For $0\le x<1$, $x^n\to0$, so the limit is $x$.
At $x=1$, $f_n(1)=0$ for every $n$. Hence
\[
 f(x)=\begin{cases}x&0\le x<1,\\0&x=1.\end{cases}
\]
\step{(b) Apply the continuity obstruction.}
Each $f_n$ is a polynomial and therefore continuous on $[0,1]$.
The limit is discontinuous at $1$, because its left-hand limit is $1$
and its value there is zero. By Theorem~24.3 the convergence is not uniform.
\answer{There is pointwise convergence to the displayed function, but no
uniform convergence on $[0,1]$.}

\exercise{24.8}{Partial geometric sums at their divergent endpoint}
\goal Analyze $f_n(x)=\sum_{k=0}^n x^k$ on $[0,1]$.
\step{(a) Evaluate the finite sum where its formula is valid.}
For $0\le x<1$,
\[
 f_n(x)=\frac{1-x^{n+1}}{1-x}\longrightarrow\frac1{1-x}.
\]
At $x=1$ there are $n+1$ summands equal to $1$, so $f_n(1)=n+1\to+\infty$.
Thus there is \emph{no real-valued pointwise limit on all of $[0,1]$}.
An infinite limit at a point is not a real-valued limit function in Definition~24.1.
\step{(b) Uniform convergence would imply pointwise convergence.}
Since real-valued pointwise convergence already fails at $1$, uniform
convergence on $[0,1]$ is impossible.
\answer{Convergence occurs pointwise only on $[0,1)$, to $1/(1-x)$;
it does not occur pointwise, or uniformly, on the full requested interval.}

\exercise{24.9}{A moving peak whose integral vanishes}
\goal Study $f_n(x)=nx^n(1-x)$ on $[0,1]$.
\step{(a) Fix an input.}
At $x=0$ and $x=1$ the function is zero.
For $0<x<1$, the ratio of successive terms of $nx^n$ is
$x(1+1/n)\to x<1$, so $nx^n\to0$. Multiplication by the fixed number
$1-x$ gives $f_n(x)\to0$. Thus $f=0$ everywhere on $[0,1]$.
\step{(b) Find a peak that does not shrink to zero.}
Differentiating gives
\[
 f_n'(x)=nx^{n-1}\bigl(n-(n+1)x\bigr).
\]
The sign changes from positive to negative at $x_n=n/(n+1)$, the maximum point.
Its value is
\[
 f_n(x_n)=\frac{n}{n+1}\left(\frac{n}{n+1}\right)^n
 =\left(\frac{n}{n+1}\right)^{n+1}\longrightarrow e^{-1}>0.
\]
The last limit follows from $(1+1/n)^n\to e$.
Because the maximum error does not approach zero, convergence is not uniform.
\step{(c) Compute the integrals directly.}
Despite the preceding failure,
\[
 \int_0^1f_n(x)\,dx
 =n\left(\frac1{n+1}-\frac1{n+2}\right)
 =\frac{n}{(n+1)(n+2)}\longrightarrow0
 =\int_0^1f(x)\,dx.
\]
\answer{The limit is zero, convergence is not uniform, and the integrals
nevertheless converge to the integral of the limit. Uniform convergence is a
sufficient condition for this interchange, not a necessary one.}

\exercise{24.10}{Adding uniform limits, and the danger with products}
\goal Prove the sum rule and examine its product analogue.
\step{(a) Split the desired error into two smaller errors.}
Let $\varepsilon>0$. Uniform convergence supplies $N_1,N_2$ such that
$|f_n(x)-f(x)|<\varepsilon/2$ for $n>N_1$ and
$|g_n(x)-g(x)|<\varepsilon/2$ for $n>N_2$, in each case for all $x\in S$.
For $n>\max\{N_1,N_2\}$ the triangle inequality gives
\[
 |(f_n+g_n)(x)-(f+g)(x)|
 \le |f_n(x)-f(x)|+|g_n(x)-g(x)|<\varepsilon.
\]
This is uniform convergence of the sums.
\step{(b) Test the product claim before accepting it.}
Products need not converge uniformly: on $S=\R$, let $f_n(x)=x$ and
$g_n(x)=1/n$. Their limits are $f(x)=x$ and $g(x)=0$, both uniform.
But $f_n(x)g_n(x)=x/n$ has value $1$ at $x=n$, so it fails to converge
uniformly to $fg=0$. A small error multiplied by an unbounded function
can become large; this is the missing issue in a naive product proof.
\answer{Sums preserve uniform convergence; unrestricted products do not.}

\exercise{24.11}{The product counterexample in full}
\goal Verify every part of the example $f_n(x)=x$, $g_n(x)=1/n$ on $\R$.
\step{(a) Compute the two errors separately.}
With $f(x)=x$ we have $|f_n(x)-f(x)|=0$ for every $n,x$.
With $g(x)=0$ we have $|g_n(x)-g(x)|=1/n$, independently of $x$.
The first error is already zero, and the second tends to zero, so both
convergences are uniform.
\step{(b) Compute the product error at a suitable input.}
For each $n$, use $x_n=n$:
\[
 |f_n(x_n)g_n(x_n)-f(x_n)g(x_n)|=|n/n-0|=1.
\]
No common index can make all these errors smaller than $1/2$.
\answer{$f_n\to f$ and $g_n\to g$ uniformly, but $f_ng_n\not\to fg$
uniformly on $\R$.}

\exercise{24.12}{The supremum-error characterization}
\goal Prove that $f_n\to f$ uniformly on $S$ if and only if
$E_n=\sup_{x\in S}|f_n(x)-f(x)|\to0$.
For an empty domain the assertion is vacuous; take its maximum error to be zero.
For a nonempty domain, $E_n$ is initially allowed to be $+\infty$.
\step{Uniform convergence implies vanishing suprema.}
Given $\varepsilon>0$, use uniform convergence with $\varepsilon/2$.
For sufficiently large $n$, all the errors are smaller than $\varepsilon/2$,
so their supremum satisfies $0\le E_n\le\varepsilon/2<\varepsilon$.
Notice the non-strict inequality at the supremum: a least upper bound need
not be attained. The extra factor $1/2$ avoids any difficulty.
\step{Vanishing suprema imply uniform convergence.}
Given $\varepsilon>0$, choose $N$ with $E_n<\varepsilon$ for all $n>N$.
Every individual error is at most its supremum, so
$|f_n(x)-f(x)|\le E_n<\varepsilon$ for all $x\in S$ and $n>N$.
This is exactly the definition of uniform convergence.
\answer{The two conditions are equivalent. The supremum is over the input
variable first; the limit is then taken over the index.}

\exercise{24.13}{Uniform limits of uniformly continuous functions}
\goal Show that a uniform limit $f$ of uniformly continuous $f_n$ on $(a,b)$
is uniformly continuous there.
\idea Approximate $f$ by one fixed $f_n$, and let that function control
nearby inputs. There will be three errors, each allocated $\varepsilon/3$.
\step{Choose one approximating function before choosing the inputs.}
Given $\varepsilon>0$, uniform convergence supplies an index $m$ with
$|f_m(t)-f(t)|<\varepsilon/3$ for all $t\in(a,b)$.
Uniform continuity of this particular $f_m$ supplies $\delta>0$ such that
$|f_m(x)-f_m(y)|<\varepsilon/3$ whenever $x,y\in(a,b)$ and $|x-y|<\delta$.
\step{Join the three estimates.}
For such $x,y$,
\[
 |f(x)-f(y)|\le |f(x)-f_m(x)|+|f_m(x)-f_m(y)|+|f_m(y)-f(y)|
 <\varepsilon.
\]
Neither $m$ nor $\delta$ was chosen using $x$ or $y$, so this is uniform
continuity, rather than merely continuity at a single point.
\answer{The uniform limit is uniformly continuous. The same proof works
on any subset of $\R$, not just an interval.}

\exercise{24.14}{A narrow peak near zero}
\goal Analyze $f_n(x)=nx/(1+n^2x^2)$ and the proposed limit $f=0$.
\step{(a) Fix $x$ and distinguish zero from nonzero.}
At $x=0$ every term is zero. If $x\ne0$, dividing by $n^2x^2$ gives
\[
 f_n(x)=\frac{1/(nx)}{1+1/(n^2x^2)}\longrightarrow0.
\]
Thus pointwise convergence holds on $\R$.
\step{(b) Test the shrinking input $x_n=1/n$ on $[0,1]$.}
At this input, $f_n(x_n)=1/2$. Therefore uniform convergence to zero fails
on $[0,1]$. Each fixed positive input eventually lies past the peak, but
the peak itself remains of height $1/2$.
\step{(c) Keep a fixed positive distance from zero.}
For $x\ge1$,
\[
 0\le\frac{nx}{1+n^2x^2}\le\frac{nx}{n^2x^2}
 =\frac1{nx}\le\frac1n.
\]
The upper bound tends to zero independently of $x$, proving uniform
convergence on $[1,\infty)$.
\answer{Pointwise convergence holds everywhere; uniform convergence fails
on $[0,1]$ and holds on $[1,\infty)$.}

\exercise{24.15}{A boundary layer with limiting height one}
\goal Analyze $f_n(x)=nx/(1+nx)$ on $[0,\infty)$.
\step{(a) Identify the limit.}
At zero the value is zero. If $x>0$, then
$f_n(x)=1-1/(1+nx)\to1$. Thus $f(0)=0$ and $f(x)=1$ for $x>0$.
\step{(b) Inspect the interval containing the boundary layer.}
The continuous $f_n$ have a discontinuous pointwise limit at zero, so
convergence cannot be uniform on $[0,1]$.
Explicitly, $x_n=1/n$ gives $f_n(x_n)=1/2$ and $f(x_n)=1$, an error of $1/2$.
\step{(c) Bound the error on $[1,\infty)$.}
For $x\ge1$ the limit is $1$, and
\[
 |f_n(x)-1|=\frac1{1+nx}\le\frac1{1+n}\longrightarrow0.
\]
This is an input-independent bound and hence proves uniform convergence.
\answer{The limit jumps from $0$ to $1$ at zero; convergence is not uniform
on $[0,1]$ but is uniform on $[1,\infty)$.}

\exercise{24.16}{A boundary layer with an unbounded limiting function}
\goal Analyze $f_n(x)=nx/(1+nx^2)$ for $x\ge0$.
\step{(a) Compute the pointwise limit.}
At zero the value is zero. For $x>0$, division by $n$ gives
$f_n(x)=x/(1/n+x^2)\to1/x$.
Thus $f(0)=0$ and $f(x)=1/x$ for $x>0$.
\step{(b) Show failure of uniform convergence on $[0,1]$.}
One obstruction is the discontinuity, indeed unboundedness, of the limit near
zero. More directly, at $x_n=1/\sqrt n$,
\[
 f_n(x_n)=\frac{\sqrt n}{2},\qquad f(x_n)=\sqrt n,
 \qquad |f_n(x_n)-f(x_n)|=\frac{\sqrt n}{2}.
\]
These errors grow, so they certainly do not tend uniformly to zero.
\step{(c) Estimate the error away from the singular endpoint.}
For $x\ge1$,
\[
 \left|\frac{nx}{1+nx^2}-\frac1x\right|
 =\frac1{x(1+nx^2)}\le\frac1{1+n}\longrightarrow0.
\]
\answer{The displayed piecewise $f$ is the limit; convergence fails to be
uniform on $[0,1]$ and is uniform on $[1,\infty)$.}

\exercise{24.17}{Passing to the limit while the input also moves}
\goal If $f_n\to f$ uniformly on $[a,b]$, the $f_n$ are continuous,
and $x_n\to x$ in that interval, prove $f_n(x_n)\to f(x)$.
\step{First secure continuity of the limiting function.}
By Theorem~24.3, the uniform limit $f$ is continuous on $[a,b]$.
Also $x\in[a,b]$, since this interval is closed. Thus continuity gives
$f(x_n)\to f(x)$.
\step{Separate the changing function from the changing input.}
For every $n$,
\[
 |f_n(x_n)-f(x)|\le |f_n(x_n)-f(x_n)|+|f(x_n)-f(x)|.
\]
Given $\varepsilon>0$, uniform convergence makes the first term smaller
than $\varepsilon/2$ for all sufficiently large $n$, \emph{even at the
moving input $x_n$}. Continuity and $x_n\to x$ make the second term smaller
than $\varepsilon/2$ for all sufficiently large $n$.
Taking $n$ beyond both thresholds proves the desired limit.
\answer{$f_n(x_n)\to f(x)$. Uniform convergence is what allows us to
substitute a different input at every index.}
